\section{Discussion}
\label{sec:discussion}

\subsection{Change detection is not source identification}

The central methodological distinction in this study is between detecting temporal change and identifying a neuron-specific source. A two-frame operator can respond strongly at the onset of a burst while discarding slow fluorescence morphology or preserving activity that is equally visible in surrounding tissue. The near-equivalence between the learned ICA component and signed difference is therefore not merely a negative result. It identifies the operation that the learning procedure selected under this data geometry and prevents an unsupported interpretation of the component as an independently recovered neuronal signal.

This distinction is consistent with the evolution from PCA/ICA cellular-signal
sorting to modern modular calcium-imaging methods, which separately model
spatial footprints, temporal traces, background activity, and deconvolved
events \citep{mukamel2009automated,pnevmatikakis2016simultaneous,giovannucci2019caiman,zhou2018cnmfe}.
The appropriate comparison for NeuRev is not whether every representation
resembles one of these packages, but whether each proposed operator contributes
a measurable function: temporal-change sensitivity, spatial localization,
nuisance removal, waveform preservation, or candidate ranking.

\subsection{A shared burst waveform can coexist with local neuronal information}

The current evidence supports a strong event-aligned Raw waveform while showing that nearby tissue contains a similar morphology. This pattern is compatible with several mechanisms, including widespread neural recruitment, neuropil contamination, motion-coupled fluorescence, common illumination or acquisition dynamics, or spatially broad indicator responses. The data do not yet identify their relative contributions. Local annulus subtraction removes part of the shared component and leaves event-associated amplitude, but it can also suppress real spatially broad neural signal. Consequently, residualization should be evaluated as a measurement transformation, not assumed to yield ground truth.

The proposed functional and tensor decompositions make this ambiguity explicit. A dominant rank-one structure would support a shared temporal waveform modulated by site and burst gains. Reproducible additional components could indicate kinetic or spatial substructure, but they would still require matched spatial controls and held-out stability before being interpreted as neuron-specific. Conversely, failure to find stable additional components would bound the amount of neuron-specific waveform information identifiable from this recording.

\subsection{Feature value depends on where the search problem begins}

The new complete-trace panel and the earlier protected full-field screen address
different stages of the pipeline. Once a confirmed center is supplied, Raw and
carrier amplitude already place nearly every annotated event above every
same-duration quiet window, leaving little headroom for a more elaborate
transform. When the location is unknown and the candidate budget is scarce,
local coherence improves known-positive recovery by requiring neighboring
pixels to express compatible dynamics. This is consistent with correlation and
peak-to-noise seeding in CNMF-E and with the broader decomposition of spatial,
temporal, and background terms in CaImAn
\citep{zhou2018cnmfe,giovannucci2019caiman}.

This division argues for a compact role-specific stack. Amplitude should remain
the trace-level carrier; coherence should supply spatial proposal/ranking
evidence; lagged recurrence should remain secondary temporal context and must
not be interpreted causally. Annular subtraction and variance stabilization
remain valuable diagnostics, but local neuropil correction can either reduce
contamination or remove genuine shared activity and therefore requires
dataset-specific validation \citep{keemink2018fissa}. Likewise, the exponential
matched-filter comparator is not a substitute for generative spike inference.
Methods such as OASIS model calcium kinetics explicitly, while benchmarking
shows that performance depends on indicator, sampling, preprocessing, task
metric, and electrophysiological ground truth
\citep{friedrich2017oasis,berens2018benchmarking}.

The consensus and multiscale-persistence extensions are most interesting as
observability features. They do not improve event localization over the carrier,
but they retain measurable cross-burst site-rank variation after amplitude
percentiles saturate. That stability is hypothesis-generating evidence for a
site reliability profile, not confirmation of a neuronal phenotype. Because
the same recording motivated and evaluated the extensions, their definitions
must now be frozen and transferred without tuning to an independent recording.

The automated challenge suite narrows these roles further. Raw amplitude,
carrier, coherence, and recurrence remain strong under frame-level retrieval,
small timing shifts, and moderate coordinate displacement. A richer grouped
recovery model improves discrimination in point estimates, but the
site-bootstrap interval crosses zero; with only 12 primary misses, this is not
confirmatory feature complementarity. The high repeatability but weak retrieval
of annular contrast is an equally important warning: a stable statistic may
encode persistent local contamination or acquisition geometry rather than
useful event signal.

The subsequent deep dives make the engineering roles more concrete. Coherence
and lag produced much larger early-budget advantages in their native proposal
streams than when restricted to a common proposal universe, which argues
against treating them as rerankers alone. Lag added more repeated-split recovery
value than coherence, whereas the full role-specific stack remained the best
point estimate. A one-frame-rise, 20-frame-decay filter was selected in every
leave-one-burst-out fold, supporting a small kinetic bank as a transparent
alternative to a single arbitrary time constant. None of these results removes
the need for new recordings or bounded-field truth: only 12 all-lane misses are
available, and the proposal contrast is descriptive rather than causal.

The canonical-v8 identity-aware audit further shows why localization and
identity resolution cannot be collapsed into a single response-maximization
step. Eight of the 12 local consensus shifts entered the neighborhood of an
existing labeled geometry, including six shifts toward a different canonical
identity. The stronger shifted waveform may therefore belong to a neighboring
source rather than reveal a corrected center for the missed observation. This
supports joint footprint or assignment-aware localization and makes automatic
recentring an unsafe repair rule in crowded regions. The four identity-clear
cases remain useful targeted hypotheses, but noise, ranking, and
non-maximum-suppression still require separate adjudication.

The highest-value feature work is consequently identity-conditional. First,
the binary recovered-versus-missed outcome should be decomposed into recovered,
identity-clear miss, and identity-ambiguous collision. With only four
identity-clear misses, this is initially a descriptive or penalized model, not
a new confirmatory classifier. Second, spatial-offset vectors should be tested
for agreement across Raw, ICA, local standardization, and repeated bursts; a
stable direction supports localization error, whereas stage- or burst-specific
motion suggests contamination or acquisition drift. Third, point sampling
should be compared with fixed disks and footprint-weighted extraction to test
whether crescent, elliptical, or multi-lobed illumination explains fragile
center traces.

Fourth, crowded-field features should quantify the target-to-neighbor peak
ratio, nearest competing maximum, saddle depth, NMS suppression margin, and
local response-field topology. Fifth, original and candidate traces should be
regressed jointly on neighboring labeled traces. Large neighbor partial
correlation or a better nonnegative mixture fit would indicate source leakage
rather than successful recentering. Sixth, footprint morphology---including
eccentricity, solidity, convexity deficit, boundary sharpness, radial
asymmetry, and multi-lobe evidence---should be evaluated for repeatability and
incremental failure stratification. These tests use the existing movies and do
not require human labels for every candidate, but biological identity and
precision still require bounded review.

The expanded automated suite provides a stopping rule for incremental feature
engineering. The full panel clearly exceeds carrier-only internal prediction,
but its most stable feature is variance-stabilized annular context, which likely
combines observability, background, and acquisition geometry. Multiscale
persistence and lag recurrence add complementary context, whereas a simple
distance graph does not explain recovery. The next scientific gain is therefore
unlikely to come from another small transform on the same bursts. It requires a
change in evidence regime: independent-recording transfer, bounded
identity-complete truth, or a realistic movie generator with exact source
identity.

These results match the modular logic of established calcium-imaging methods,
which distinguish spatial filtering or footprints from temporal calcium
dynamics and background terms
\citep{vogelstein2010deconvolution,pnevmatikakis2016simultaneous,giovannucci2019caiman}.
They do not show that the present compact operators recover spikes or demixed
sources. Realistic simulation frameworks add anatomy, optical propagation,
motion, indicator kinetics, and sensor noise and can expose false positives
that strong extracted time courses alone do not reveal \citep{song2021naomi}.
The monotonic one-dimensional injection result therefore justifies retaining
the matched filter for transfer testing, while also defining why movie-level
simulation is the next synthetic step.

\subsection{Stable observability may be the most productive neuron-level target}

Detailed activation morphology need not be the only reproducible property of a neuron. A site may repeatedly exhibit high amplitude, strong SNR, compact spatial support, weak annulus coupling, or reliable candidate ranking even when its exact waveform varies. We refer to this profile as a measurement or observability phenotype. The term is deliberately nonbiological: stable visibility can arise from expression level, optical depth, focus, cell size, local crowding, acquisition field, or true response magnitude.

Quantifying observability has direct methodological value. It can explain why some known neurons are systematically recovered while others are repeatedly missed, reveal whether apparent algorithm failures are actually proposal or acquisition failures, and guide the next annotation or data-collection effort. A detector that improves only for already bright, isolated neurons is less scientifically useful than one that expands coverage of difficult but credible sites. The proposed variance-component and recoverability models therefore treat site-associated heterogeneity as a primary result rather than nuisance variation.

\subsection{Incomplete annotation changes what detection metrics mean}

Sparse positive labels make recall of known observations identifiable but leave full-field precision unidentified. Treating every unmatched candidate as background would penalize the algorithm for discovering neurons omitted by the annotation process and would favor conservative methods that reproduce the labeling bias. Positive--unlabeled learning theory emphasizes that conclusions depend on how positives enter the labeled set, particularly when labeling is not random \citep{bekker2020pu}.

The bounded-field annotation study is designed to convert one small region from positive--unlabeled to exhaustively reviewed data. Its value is disproportionate to its size: it permits a local precision estimate, a false-candidate taxonomy, an estimate of hidden-positive prevalence, and an assessment of whether candidate score calibrates review confidence. The rest of the full field remains useful for known-positive recovery and discovery, but it should not be assigned a fictitious negative class.

\subsection{Implications for unsupervised learning in NeuRev}

Avoiding large supervised models is appropriate given the current data volume, but unsupervised learning does not remove the need for scientific evaluation. An unsupervised representation can learn acquisition structure, global burst timing, local covariance, or a simple derivative. The relevant question is whether it exposes a stable and useful statistic that an interpretable comparator does not already provide.

The confirmatory program therefore places learned and analytic operators in the same frozen panel. Pairwise ICA remains as an equivalence control; local coherence and lagged recurrence are tested for ranking utility; residualization is tested for spatial specificity and fidelity; and Raw remains the carrier against which information loss is measured. New learning is warranted only when the existing analysis identifies a missing operation or a reproducible residual structure that cannot be represented by the compact panel.

For subsequent max-envelope studies, we distinguish
\emph{causal local standardized surprise} (the existing LS/MSLN output) from
\emph{positive instantaneous evidence} at the current frame and coordinate,
its \emph{contextual upper envelope} over a declared causal support, and their
\emph{envelope agreement}. The resulting product is termed
\emph{envelope-gated evidence}; multiplying instantaneous evidence by envelope
agreement gives \emph{agreement-attenuated evidence}. This vocabulary is
operator-specific and deliberately does not rename the established
\emph{local coherence} feature, which remains a causal pixel-to-neighborhood
correlation. The bounded pilot defining these terms is non-claim-bearing and
does not establish detection utility. In the subsequent all-occurrence
comparison, temporal agreement attenuation was effectively indistinguishable
from amplitude squaring under peak-localization ranking, whereas adding a
3-by-3 spatial envelope reduced known-center localization, especially after
temporal-envelope-to-LS-to-ICA processing. This follows from the operator and
estimand: at a local maximum included in its causal envelope, the upper
envelope equals the instantaneous evidence and the attenuation vanishes.
Temporal agreement attenuation should therefore remain a morphology-sensitive
diagnostic unless it succeeds under integrated-energy, duration, or bounded
full-field specificity endpoints.

A subsequent visual-first morphology audit confirms that this caution should
not be mistaken for rejection of temporal pooling. Across all 106 confirmed
occurrences, the five-frame upper envelope broadened and prolonged event
evidence, whereas agreement-attenuated evidence preserved the aligned peak,
increased core concentration, and reduced the post-peak shoulder in Raw and
both learned representations. Envelope-gated evidence produced stronger event
concentration but behaves partly as an amplitude-emphasizing transform. These
operators therefore have different potential product roles: short causal
memory, proposal reinforcement, and peak-preserving temporal sharpening. They
remain within-recording measurement features; whether any improves candidate
specificity requires bounded-field truth, and waveform analyses must preserve
the unattenuated carrier because attenuation deliberately changes decay shape.

The spatial-ICA results sharpen this conclusion. Individual filters varied too much across seeds and held-burst exclusions to support component-wise anatomical interpretation, yet dense reconstruction was nearly invariant and preserved event-map morphology, peak timing, and known-positive recall. Basis instability and representation stability are not contradictory: multiple rotated or redistributed bases can span a similar useful signal subspace. For annotation, spatial ICA can remain a useful visual and quantitative representation without implying that any displayed component is a uniquely identifiable neuron.

Objective translated-control comparisons support localized morphology at labeled coordinates, but they do not turn sparse positives into exhaustive truth. Likewise, the class analyses provide plausible structure---especially the candidate-assisted boundary-sharpness gradient---without corrected evidence that the classes encode spatial-ICA morphology or reviewer certainty. The frozen classes should remain measurement-event archetypes rather than neuron types.

\subsection{Limitations}

The study has several non-negotiable limitations. First, one recording cannot establish generalization across animals, preparations, sessions, or microscopes. Second, four shared bursts provide limited independent temporal replication and make burst-level variance estimates imprecise. Third, current labels are sparse and partly uncertain, and the same bursts have informed multiple exploratory decisions. Fourth, pre-repair identity-sensitive numbers remain provisional even though later analyses preserve immutable site geometry. Fifth, annulus subtraction is not a pure nuisance-removal operation. Sixth, candidate-based metrics outside the bounded field cannot identify precision or specificity. Seventh, any coactivity or lag analysis is descriptive because shared drive and calcium kinetics can generate apparent network structure. Eighth, the spatial-ICA fit is transductive within one recording, its positive encoding clips negative reconstructed values, and its manual review panel has been deferred; objective morphology and reconstruction stability do not establish reviewer utility or biological source separation. Ninth, the complete-trace panel conditions on confirmed centers and reuses previously analyzed bursts; its ceiling-limited event-localization metric neither validates the two new features nor ranks their full-field specificity. Tenth, the automated recovery model has only 12 primary misses and its improvement interval crosses zero. Eleventh, circular shifts and simplified trace injections validate alignment and operator behavior but do not reproduce anatomy, optics, motion, or biological ground truth. Twelfth, the canonical-v8 recentering candidates are response-selected within the same recording, and the identity guard uses sparse labeled geometry rather than exhaustive biological source truth.

\begin{table}[!htbp]
\centering
\caption{Limitations and the corresponding manuscript safeguard.}
\label{tab:limitations}
\small
\begin{tabularx}{\linewidth}{>{\raggedright\arraybackslash}p{0.30\linewidth} X}
\toprule
Limitation & Safeguard \\
\midrule
One canonical recording & Restrict claims to within-recording measurement structure; require an independent recording for external generalization. \\
Four shared burst intervals & Use leave-one-burst-out and burst-exclusion sensitivity; avoid asymptotic burst-level claims. \\
Sparse, partly uncertain positives & Use positive--unlabeled metrics and preserve confidence/provenance; estimate precision only in a bounded exhaustive field. \\
Identity/geometry collision in existing assays & Repair before trace aggregation; preserve original sites and regenerate all identity-sensitive results. \\
Response-selected recentering can transfer identity & Require nearest-label and footprint-assignment guards; never count a stronger shifted trace as recovery by itself. \\
Repeated exploratory use of the same bursts & Freeze the compact panel and estimands before a protected confirmatory rerun. \\
Local annulus may contain true neural signal & Report Raw, annulus and residual jointly; treat subtraction as a transformation, not ground truth. \\
Acquisition heterogeneity and saturation & Characterize noise, field effects, motion and saturation before interpreting stable observability. \\
Shared drive can induce network correlations & Keep connectivity analyses descriptive and condition on global/local common mode. \\
\bottomrule
\end{tabularx}
\end{table}


These limitations do not make the dataset scientifically unusable. They determine the level at which it should be analyzed. The strongest contribution is a transparent within-recording measurement study, accompanied by explicit criteria for the additional data and annotation required to promote each claim.

\subsection{Highest-value next evidence}

The most valuable immediate addition is not a wider hyperparameter sweep. It
is exhaustive review of one bounded field using the frozen carrier, coherence,
and recurrence panel, so proposal coverage, ranking, and precision can be
separated. The most valuable future dataset is an independently acquired
recording that satisfies a frozen quality and annotation contract. Such a
recording would test the shared-versus-specific decomposition and determine
whether the consensus and multiscale-persistence site rankings transfer without
retuning. For synthetic work, an anatomy- and optics-aware movie simulation is
more informative than further widening the present one-dimensional injection
grid.
